\begin{recipe}{Gevulde varkensrollade met pesto en roomkaas}
\kicker[Hoofdgerecht,Vlees,Italiaans]{Hoofdgerecht · vlees · Italiaans}
\index[register]{Gevulde varkensrollade met pesto en roomkaas}
\index[register]{Varkensvlees}
\index[register]{Pesto}
\index[register]{Roomkaas}
\index[register]{Porchetta}

\meta{30 minuten + 75 minuten oven \dvd Voor 6 personen}

\lettrine{E}{en} gevulde varkensrollade in de stijl van een porchetta, met een
mooie spiraal van roomkaas, pesto, spinazie en geroosterde paprika. Al meermaals
gemaakt voor kerst en voor bezoek. Lekker met gebakken aardappels.

\blockrule{Ingrediënten}
\begin{ingredients}
  \ing{ca.\ 1,5 kg}{varkenshaas}\index[register]{Varkensvlees}
  \ing{125 g}{roomkaas, zacht}\index[register]{Roomkaas}
  \ing{400 g}{geroosterde rode paprika, uitgelekt (1 pot, gebruik wat past)}\index[register]{Paprika}
  \ing{100 g}{spek, ongebakken}
  \ing{8 plakken}{koude prosciutto}
  \ing{90 g}{verse pesto}\index[register]{Pesto}
  \ing{200 g}{bladspinazie}
  \ingb{olijfolie}
  \ingb{paprikapoeder, zout, versgemalen peper}
\end{ingredients}

\blockrule{Bereiding}
\tip{Gebruik een vleesthermometer. Bij 63\,°C kerntemperatuur is het vlees gaar
en nog sappig. Laat de rollade altijd een paar minuten rusten voor je hem
aansnijdt.}
\begin{steps}
  \item Leg het vlees plat op een snijplank. Snijd het open als een boek
        (vlindermethode) zodat je een groot plat stuk krijgt. Dek af met
        vershoudfolie en sla gelijkmatig plat met een vleeshamer tot ca.\ 1 cm
        dikte.
  \item Verdeel de spinazie, geroosterde paprika, pesto, spek, prosciutto en
        roomkaas gelijkmatig over het vlees. Breng op smaak met zout en peper.
  \item Rol het vlees strak op vanaf de lange kant, zoals een opgerolde handdoek.
        Stop de uiteinden in en bind het geheel vast met keukentouw. Wrijf de
        buitenkant in met olijfolie en bestrooi met paprikapoeder, zout en peper.
  \item Verwarm de oven voor op 175\,°C (boven- en onderwarmte). Leg de rollade in een braadpan en bak
        60--90 minuten, tot de kerntemperatuur 63\,°C bereikt.
  \item Laat de rollade een paar minuten rusten en snijd hem daarna in plakken.
\end{steps}
\end{recipe}
